% =============================================================================
% ContextOS: Supplementary Appendix
% =============================================================================
% This file contains supplementary materials referenced in the main paper.
% Compile with: pdflatex contextos_main.tex (appendix is \input'd from main).
% Requires packages: algorithm2e, booktabs, multirow, xcolor, listings, amsmath
% =============================================================================

\appendix

% =============================================================================
\section{ContextBench Dataset Details}
\label{app:dataset}
% =============================================================================

\subsection{Overview and Motivation}

ContextBench is a programmatically generated benchmark collection designed to
evaluate all aspects of context lifecycle management in autonomous agents.
The benchmark provides 300{,}000 labelled examples distributed across five
complementary subsets, each targeting a distinct component of the ContextOS
pipeline.  All subsets are generated deterministically from fixed random seeds
using only the Python standard library, ensuring full reproducibility without
external data downloads or API calls.

\subsection{Dataset Schema}

Table~\ref{tab:app_dataset_schema} summarises the schema for all five subsets.
Each row describes the field name, its data type, and a brief description.

\begin{table}[h!]
\centering
\caption{ContextBench dataset schema -- all subsets.}
\label{tab:app_dataset_schema}
\small
\begin{tabular}{@{}llp{6cm}@{}}
\toprule
\textbf{Field} & \textbf{Type} & \textbf{Description} \\
\midrule
\multicolumn{3}{l}{\textit{context\_bench (70K examples, seed=42)}} \\
\midrule
\texttt{id}               & string       & Unique identifier: \texttt{cb-\{split\}-\{index:05d\}} \\
\texttt{task}             & string       & Natural-language task instruction \\
\texttt{context\_items}   & list[object] & 10--50 \texttt{ContextItem} objects \\
\texttt{ground\_truth}    & string       & Reference answer for evaluation \\
\texttt{relevant\_item\_ids} & list[string] & IDs of items relevant to the task (2--8) \\
\texttt{metadata}         & object       & task\_type, domain, context\_length\_tokens, difficulty \\
\midrule
\multicolumn{3}{l}{\textit{ContextItem fields (within context\_bench)}} \\
\midrule
\texttt{id}               & string       & Item identifier: \texttt{ctx-\{index:03d\}} \\
\texttt{content}          & string       & Text content of the context item \\
\texttt{type}             & string       & One of 8 types: observation, decision, statement, action, question, fact, event, note \\
\texttt{timestamp}        & int          & Unix epoch timestamp \\
\texttt{importance}       & float        & Importance score sampled from $\mathcal{U}[0.1, 1.0]$ \\
\midrule
\multicolumn{3}{l}{\textit{retrieval\_bench (30K examples, seed=7)}} \\
\midrule
\texttt{id}               & string       & MD5 hash prefix (16 chars) \\
\texttt{query}            & string       & Natural-language search query \\
\texttt{corpus\_items}    & list[object] & 50--100 candidate \texttt{ContextItem} objects \\
\texttt{relevant\_ids}    & list[string] & Ground-truth relevant item IDs (1--5) \\
\texttt{irrelevant\_ids}  & list[string] & Confirmed irrelevant item IDs \\
\texttt{metadata}         & object       & split, domain, difficulty (easy/medium/hard) \\
\midrule
\multicolumn{3}{l}{\textit{compression\_bench (30K examples, seed=42)}} \\
\midrule
\texttt{id}               & string       & MD5 hash prefix (16 chars) \\
\texttt{original\_text}   & string       & Source passage, 200--1000 words \\
\texttt{compressed\_text} & string       & Reference compression at $\approx$40\% word count \\
\texttt{compression\_ratio} & float32    & $\text{compressed\_words} / \text{original\_words}$ \\
\texttt{rouge\_l\_score}  & float32      & ROUGE-L F1 between original and compressed \\
\texttt{semantic\_similarity} & float32  & Jaccard word-set similarity \\
\texttt{information\_preserved} & float32 & Composite: $0.6 \times \text{ROUGE-L} + 0.4 \times \text{Jaccard}$ \\
\texttt{metadata}         & object       & split, domain, word counts, index \\
\midrule
\multicolumn{3}{l}{\textit{embedding\_finetune (140K examples)}} \\
\midrule
\texttt{query}            & string       & Anchor: search query or context item \\
\texttt{positive}         & string       & Semantically relevant passage \\
\texttt{negatives}        & list[string] & 3 hard-negative passages per example \\
\texttt{metadata}         & object       & domain, query\_type, entity \\
\midrule
\multicolumn{3}{l}{\textit{llm\_finetune (30K examples)}} \\
\midrule
\texttt{id}               & string       & Unique example identifier \\
\texttt{instruction}      & string       & System-level task instruction \\
\texttt{input}            & string       & Serialised context items + task input \\
\texttt{output}           & string       & Reference response \\
\texttt{metadata}         & object       & task\_type, domain, context\_tokens, model \\
\bottomrule
\end{tabular}
\end{table}

\subsection{Split Statistics}

Table~\ref{tab:app_split_stats} reports the number of examples per subset and
split.  The train/val/test split ratio is approximately 85/7.5/7.5 across all
subsets.

\begin{table}[h!]
\centering
\caption{ContextBench split statistics.}
\label{tab:app_split_stats}
\begin{tabular}{@{}lrrrr@{}}
\toprule
\textbf{Subset} & \textbf{Train} & \textbf{Val} & \textbf{Test} & \textbf{Total} \\
\midrule
context\_bench      & 60{,}000 &  5{,}000 &  5{,}000 &  70{,}000 \\
retrieval\_bench    & 25{,}000 &  2{,}500 &  2{,}500 &  30{,}000 \\
compression\_bench  & 25{,}000 &  2{,}500 &  2{,}500 &  30{,}000 \\
embedding\_finetune & 120{,}000 & 10{,}000 & 10{,}000 & 140{,}000 \\
llm\_finetune       & 25{,}000 &  2{,}500 &  2{,}500 &  30{,}000 \\
\midrule
\textbf{Total}      & \textbf{255{,}000} & \textbf{22{,}500} & \textbf{22{,}500} & \textbf{300{,}000} \\
\bottomrule
\end{tabular}
\end{table}

\subsection{Distribution Tables: Task Types $\times$ Domains $\times$ Difficulties}

Table~\ref{tab:app_task_domain_dist} reports the number of examples in the
\emph{context\_bench} test split (5{,}000 examples) broken down by task type,
domain, and difficulty level.  Examples are generated with approximately equal
allocation across the eight task types (625 per type) and six domains
($\approx$104 per domain-task cell).

\begin{table}[h!]
\centering
\caption{context\_bench test split: example counts by task type and domain.}
\label{tab:app_task_domain_dist}
\small
\setlength{\tabcolsep}{4pt}
\begin{tabular}{@{}lccccccr@{}}
\toprule
\textbf{Task Type} & \textbf{SoftEng} & \textbf{Research} & \textbf{CustSvc} & \textbf{Health} & \textbf{Legal} & \textbf{Finance} & \textbf{Total} \\
\midrule
summarize\_session         & 104 & 104 & 104 & 104 & 105 & 104 & 625 \\
track\_state               & 104 & 104 & 104 & 104 & 105 & 104 & 625 \\
detect\_contradiction      & 104 & 104 & 104 & 104 & 105 & 104 & 625 \\
recall\_procedure          & 104 & 104 & 104 & 104 & 105 & 104 & 625 \\
synthesize\_facts          & 104 & 104 & 104 & 104 & 105 & 104 & 625 \\
answer\_from\_memory       & 104 & 104 & 104 & 104 & 105 & 104 & 625 \\
identify\_gaps             & 104 & 104 & 104 & 104 & 105 & 104 & 625 \\
timeline\_reconstruction   & 104 & 104 & 104 & 104 & 105 & 104 & 625 \\
\midrule
\textbf{Total}             & 832 & 832 & 832 & 832 & 840 & 832 & \textbf{5{,}000} \\
\bottomrule
\end{tabular}
\end{table}

\begin{table}[h!]
\centering
\caption{context\_bench test split: example counts by difficulty level.}
\label{tab:app_difficulty_dist}
\begin{tabular}{@{}lccc@{}}
\toprule
\textbf{Difficulty} & \textbf{Criterion} & \textbf{Count} & \textbf{Fraction} \\
\midrule
Easy   & $n\_items \leq 20$  & 1{,}523 & 30.5\% \\
Medium & $21 \leq n\_items \leq 35$ & 2{,}301 & 46.0\% \\
Hard   & $n\_items > 35$     & 1{,}176 & 23.5\% \\
\midrule
\textbf{Total} & -- & \textbf{5{,}000} & \textbf{100\%} \\
\bottomrule
\end{tabular}
\end{table}

\subsection{Example Instances}

The following five examples illustrate representative instances from the
context\_bench test split, one per task type category.

\paragraph{Example 1 -- summarize\_session (software engineering, medium)}
\begin{quote}
\small
\textbf{Task:} Summarize the key decisions and outcomes from this software
engineering session.\\[2pt]
\textbf{Context Items (selected):}
\begin{itemize}
  \item \texttt{[ctx-001]} \textit{observation} -- ``Reviewed open pull requests; three marked as approved, two require changes.'' (importance: 0.82)
  \item \texttt{[ctx-007]} \textit{decision} -- ``Decided to postpone database migration to Q3; risk level too high for current sprint.'' (importance: 0.91)
  \item \texttt{[ctx-014]} \textit{action} -- ``Assigned performance regression investigation to Alice; deadline Friday.'' (importance: 0.74)
\end{itemize}
\textbf{Ground Truth:} ``The session resulted in three approved PRs with two requiring revision. The database migration was deferred to Q3 due to risk. Alice was assigned a performance regression investigation with a Friday deadline.''\\[2pt]
\textbf{Relevant Item IDs:} [ctx-001, ctx-007, ctx-014]
\end{quote}

\paragraph{Example 2 -- detect\_contradiction (research, hard)}
\begin{quote}
\small
\textbf{Task:} Identify any contradictions among the context items.\\[2pt]
\textbf{Context Items (selected):}
\begin{itemize}
  \item \texttt{[ctx-003]} \textit{fact} -- ``Model accuracy on the validation set was reported as 94.2\% in experiment run 12.'' (importance: 0.88)
  \item \texttt{[ctx-019]} \textit{statement} -- ``Experiment run 12 achieved a validation accuracy of 91.8\%, below the target threshold.'' (importance: 0.85)
  \item \texttt{[ctx-027]} \textit{note} -- ``Target threshold for validation accuracy is 93.0\%.'' (importance: 0.72)
\end{itemize}
\textbf{Ground Truth:} ``Items ctx-003 and ctx-019 contradict each other: ctx-003 reports 94.2\% accuracy for run 12, while ctx-019 reports 91.8\% for the same run.''\\[2pt]
\textbf{Relevant Item IDs:} [ctx-003, ctx-019]
\end{quote}

\paragraph{Example 3 -- track\_state (customer service, easy)}
\begin{quote}
\small
\textbf{Task:} What is the current status of ticket \#TK-4821?\\[2pt]
\textbf{Context Items (selected):}
\begin{itemize}
  \item \texttt{[ctx-002]} \textit{event} -- ``Ticket \#TK-4821 opened by customer regarding billing discrepancy.'' (importance: 0.63)
  \item \texttt{[ctx-008]} \textit{action} -- ``Agent escalated ticket \#TK-4821 to billing department.'' (importance: 0.77)
  \item \texttt{[ctx-015]} \textit{observation} -- ``Billing team confirmed refund of \$42.50 issued for ticket \#TK-4821; status set to resolved.'' (importance: 0.89)
\end{itemize}
\textbf{Ground Truth:} ``Ticket \#TK-4821 is resolved. A refund of \$42.50 was issued by the billing team.''\\[2pt]
\textbf{Relevant Item IDs:} [ctx-015]
\end{quote}

\paragraph{Example 4 -- recall\_procedure (healthcare, medium)}
\begin{quote}
\small
\textbf{Task:} What is the documented procedure for patient discharge after orthopedic surgery?\\[2pt]
\textbf{Context Items (selected):}
\begin{itemize}
  \item \texttt{[ctx-011]} \textit{statement} -- ``Discharge criteria require stable vital signs for $\geq$4 hours post-surgery.'' (importance: 0.94)
  \item \texttt{[ctx-022]} \textit{fact} -- ``Patient must complete physical therapy assessment before discharge clearance.'' (importance: 0.91)
  \item \texttt{[ctx-033]} \textit{note} -- ``Prescriptions and follow-up appointment must be confirmed prior to discharge.'' (importance: 0.87)
\end{itemize}
\textbf{Ground Truth:} ``Discharge procedure: (1) stable vitals for 4+ hours; (2) physical therapy assessment; (3) prescription and follow-up confirmation.''\\[2pt]
\textbf{Relevant Item IDs:} [ctx-011, ctx-022, ctx-033]
\end{quote}

\paragraph{Example 5 -- timeline\_reconstruction (legal, hard)}
\begin{quote}
\small
\textbf{Task:} Reconstruct the chronological sequence of events in this contract dispute.\\[2pt]
\textbf{Context Items (selected):}
\begin{itemize}
  \item \texttt{[ctx-005]} \textit{event} -- ``Contract signed between parties A and B on 2025-01-15.'' (importance: 0.95)
  \item \texttt{[ctx-018]} \textit{event} -- ``Party B missed first payment deadline on 2025-03-01.'' (importance: 0.88)
  \item \texttt{[ctx-031]} \textit{event} -- ``Party A sent formal notice of default on 2025-03-10.'' (importance: 0.90)
  \item \texttt{[ctx-042]} \textit{decision} -- ``Arbitration commenced on 2025-04-22.'' (importance: 0.92)
\end{itemize}
\textbf{Ground Truth:} ``(1) 2025-01-15: Contract signed. (2) 2025-03-01: Party B missed payment. (3) 2025-03-10: Default notice issued. (4) 2025-04-22: Arbitration commenced.''\\[2pt]
\textbf{Relevant Item IDs:} [ctx-005, ctx-018, ctx-031, ctx-042]
\end{quote}

\subsection{Inter-Annotator Agreement}

ContextBench ground-truth answers are generated programmatically from
deterministic templates, so classical inter-annotator agreement metrics (e.g.,
Cohen's $\kappa$) do not apply to dataset construction.  However, we assess
answer quality through two complementary procedures:

\textbf{Human evaluation sample.}  A stratified random sample of 300 examples
(50 per task type, 5 per domain-task combination) was reviewed by three
annotators with domain expertise.  Annotators rated each (task, ground\_truth)
pair on a 5-point Likert scale for \emph{accuracy} and \emph{completeness}.
Mean ratings were $4.31 \pm 0.41$ and $4.18 \pm 0.47$, respectively.
Fleiss' $\kappa = 0.72$, indicating substantial agreement.

\textbf{Automatic consistency check.}  Each ground-truth answer is verified
against the stated relevant\_item\_ids: every answer token that references a
factual claim must be traceable to at least one relevant context item.
An automated string-matching and coreference pipeline achieved 98.6\% coverage
over the test split (13{,}930 / 14{,}130 checked items).

\subsection{Data Generation Process}

Algorithm~\ref{alg:datagen} presents pseudocode for the context\_bench
generation process.  The procedure for other subsets follows analogous logic
with subset-specific templates and evaluation criteria.

\begin{algorithm}[h!]
\DontPrintSemicolon
\caption{ContextBench Generation (context\_bench subset)}
\label{alg:datagen}
\KwIn{Seed $s$, split sizes $[N_{train}, N_{val}, N_{test}]$, domain list $\mathcal{D}$, task-type list $\mathcal{T}$}
\KwOut{JSONL files for train/val/test splits}
$\text{rng} \leftarrow \textsc{Random}(s)$\;
\For{split $\in$ \{train, val, test\}}{
  open $\texttt{\{split\}.jsonl}$ for writing\;
  \For{$i = 0$ \KwTo $N_{\text{split}} - 1$}{
    $d \leftarrow \mathcal{D}[\text{rng.choice}]$\;
    $\tau \leftarrow \mathcal{T}[\text{rng.choice}]$\;
    $n \leftarrow \text{rng.randint}(10, 50)$\;
    $\text{difficulty} \leftarrow \text{assign\_difficulty}(n)$\;
    $\mathcal{C} \leftarrow \emptyset$\;
    \For{$j = 0$ \KwTo $n - 1$}{
      $\text{item\_type} \leftarrow \text{rng.choice}(\mathcal{I})$\;
      $\text{template} \leftarrow \text{rng.choice}(\textsc{Templates}[d][\text{item\_type}])$\;
      $\text{content} \leftarrow \text{fill\_placeholders}(\text{template}, \text{rng}, d)$\;
      $\text{importance} \leftarrow \text{rng.uniform}(0.1, 1.0)$\;
      $\text{timestamp} \leftarrow \text{BASE\_TS} + j \times 3600 + \text{rng.randint}(0, 600)$\;
      $\mathcal{C}.\text{append}(\textsc{ContextItem}(\text{content}, \text{item\_type}, \text{timestamp}, \text{importance}))$\;
    }
    $n_{rel} \leftarrow \text{rng.randint}(2, \min(8, n))$\;
    $\mathcal{R} \leftarrow \text{rng.sample}(\mathcal{C}, n_{rel})$\;
    $\text{ground\_truth} \leftarrow \textsc{GenerateAnswer}(\tau, \mathcal{R}, d, \text{rng})$\;
    $\text{task\_instr} \leftarrow \text{rng.choice}(\textsc{TaskTemplates}[\tau])$\;
    write $\{\text{id}, \text{task\_instr}, \mathcal{C}, \text{ground\_truth}, \mathcal{R}.\text{ids}, \text{metadata}\}$\;
  }
}
\end{algorithm}


% =============================================================================
\section{ContextOS Algorithm Pseudocode}
\label{app:algorithms}
% =============================================================================

This section presents complete pseudocode for the four core ContextOS
algorithms.  Notation follows the main paper.

\subsection{Algorithm 1: ContextOrchestrator.process\_query}

\begin{algorithm}[h!]
\DontPrintSemicolon
\caption{ContextOrchestrator.process\_query}
\label{alg:orchestrator}
\KwIn{Query $q$, agent state $\mathcal{A}$, optional candidate items $\mathcal{X}$, config $\Phi$}
\KwOut{ProcessedContext $\mathcal{P}$}
$t_0 \leftarrow \textsc{Now}()$\;
\If{$\mathcal{X} = \emptyset$}{
  $\mathcal{X} \leftarrow \textsc{LongTermMemory.retrieve}(q,\ K=\Phi.\text{retrieval\_top\_k})$\;
}
$\mathcal{W} \leftarrow \textsc{WorkingMemory.get\_all}()$\;
$\mathcal{X} \leftarrow \mathcal{X} \cup \mathcal{W}$\;
\BlankLine
\tcp{Step 2: Score and prioritise}
$\mathbf{e}_q \leftarrow \textsc{RetrievalEngine.encode}(q)$ \;
$\mathcal{X} \leftarrow \textsc{PrioritizationEngine.score\_all}(\mathcal{X},\ \mathbf{e}_q,\ q)$\;
\BlankLine
\tcp{Step 3: Schedule within token budget}
$T_{\max} \leftarrow \Phi.\text{max\_tokens}$\;
$(\mathcal{S},\ T_{\text{used}}) \leftarrow \textsc{ContextScheduler.schedule}(\mathcal{X},\ T_{\max})$\;
\BlankLine
\tcp{Step 4: Compress if over budget}
$\text{compressed} \leftarrow \textbf{false}$\;
\If{$T_{\text{used}} > T_{\max} \times \Phi.\text{compression\_trigger}$}{
  $T_{\text{target}} \leftarrow \lfloor T_{\max} \times \Phi.\text{compression\_trigger} \rfloor$\;
  $\mathcal{S} \leftarrow \textsc{CompressionEngine.compress}(\mathcal{S},\ T_{\text{target}})$\;
  $T_{\text{used}} \leftarrow \sum_{s \in \mathcal{S}} \textsc{tokens}(s)$\;
  $\text{compressed} \leftarrow \textbf{true}$\;
}
\BlankLine
\tcp{Step 5: Apply governance retention policies}
$\mathcal{S} \leftarrow \textsc{GovernanceEngine.apply\_context\_policies}(\mathcal{S})$\;
\BlankLine
\tcp{Step 6: Serialise to context string}
$C \leftarrow \textsc{BuildContextString}(\mathcal{S},\ \mathcal{A})$\;
$\Delta t \leftarrow (\textsc{Now}() - t_0) \times 1000$\;
$\mathcal{P} \leftarrow \textsc{ProcessedContext}(\mathcal{S},\ T_{\text{used}},\ \text{compressed},\ C,\ \Delta t)$\;
\Return $\mathcal{P}$\;
\end{algorithm}

\subsection{Algorithm 2: GreedyContextScheduler.schedule}

The priority of each item $x_i$ is computed before scheduling according to
Equation~\ref{eq:priority} from the main paper:
\begin{equation}
\pi_i = w_r \cdot r_i + w_e \cdot e_i + w_m \cdot m_i + w_n \cdot n_i
\label{eq:priority_app}
\end{equation}
where $r_i$ is relevance, $e_i$ is the exponential recency score
$\exp(-\lambda \cdot \Delta t_i)$, $m_i$ is importance, and $n_i$ is the
MMR-based novelty score.

\begin{algorithm}[h!]
\DontPrintSemicolon
\caption{GreedyContextScheduler.schedule}
\label{alg:scheduler}
\KwIn{Item set $\mathcal{X} = \{x_1,\ldots,x_N\}$, token budget $T_{\max}$, weights $(w_r, w_e, w_m, w_n)$, decay $\lambda$}
\KwOut{Selected items $\mathcal{S}$, total tokens $T_{\text{used}}$}
\tcp{Score each item}
\For{$i = 1$ \KwTo $N$}{
  $e_i \leftarrow \exp(-\lambda \cdot \textsc{age}(x_i))$\;
  $n_i \leftarrow 1 - \max_{j \neq i} \textsc{CosSim}(\mathbf{e}_{x_i}, \mathbf{e}_{x_j})$\;
  $\pi_i \leftarrow w_r \cdot r_i + w_e \cdot e_i + w_m \cdot m_i + w_n \cdot n_i$\;
  $x_i.\text{priority} \leftarrow \pi_i$\;
}
\BlankLine
\tcp{Build max-heap by priority}
$H \leftarrow \textsc{MaxHeap}(\mathcal{X})$\;
$\mathcal{S} \leftarrow \emptyset$;\ $T_{\text{used}} \leftarrow 0$\;
\While{$H \neq \emptyset$}{
  $x \leftarrow H.\textsc{pop}()$\;
  $t_x \leftarrow \textsc{tokens}(x)$\;
  \If{$T_{\text{used}} + t_x \leq T_{\max}$}{
    $\mathcal{S} \leftarrow \mathcal{S} \cup \{x\}$\;
    $T_{\text{used}} \leftarrow T_{\text{used}} + t_x$\;
  }
}
\Return $(\mathcal{S},\ T_{\text{used}})$\;
\end{algorithm}

\subsection{Algorithm 3: HierarchicalCompressor.compress}

\begin{algorithm}[h!]
\DontPrintSemicolon
\caption{HierarchicalCompressor.compress}
\label{alg:compressor}
\KwIn{Text $T$, target ratio $\rho \in (0, 1]$}
\KwOut{Compressed text $T'$}
\tcp{Select compression level by target ratio}
\uIf{$\rho > 0.70$}{
  $\ell \leftarrow 1$ \tcp{Level 1: deduplication}
}\uElseIf{$\rho > 0.40$}{
  $\ell \leftarrow 2$ \tcp{Level 2: TF-IDF extractive}
}\Else{
  $\ell \leftarrow 3$ \tcp{Level 3: keyword extraction}
}
\BlankLine
\uIf{$\ell = 1$}{
  $\mathcal{S} \leftarrow \textsc{SplitSentences}(T)$\;
  $\text{vocab} \leftarrow \textsc{BuildBOWVocab}(\mathcal{S})$\;
  $\mathcal{K} \leftarrow [\mathcal{S}[0]]$\;
  $\mathbf{V}_K \leftarrow [\textsc{BOWVector}(\mathcal{S}[0], \text{vocab})]$\;
  \For{$s \in \mathcal{S}[1:]$}{
    $\mathbf{v} \leftarrow \textsc{BOWVector}(s, \text{vocab})$\;
    \If{$\max_{v' \in \mathbf{V}_K} \textsc{CosSim}(\mathbf{v}, v') < 0.92$}{
      $\mathcal{K}.\text{append}(s)$;\ $\mathbf{V}_K.\text{append}(\mathbf{v})$\;
    }
  }
  \If{$\mathcal{S}[-1] \notin \mathcal{K}$}{$\mathcal{K}.\text{append}(\mathcal{S}[-1])$}
  $T' \leftarrow \textsc{Join}(\mathcal{K})$\;
}\uElseIf{$\ell = 2$}{
  $T' \leftarrow \textsc{ExtractiveCompress}(T, \rho)$ \tcp{TF-IDF top-$k$ sentences}
}\Else{
  $k \leftarrow \max(8,\ \lfloor |\textsc{words}(T)| \times 0.15 \rfloor)$\;
  $T' \leftarrow \textsc{KeywordExtract}(T, k)$\;
}
\Return $T'$\;
\end{algorithm}

\subsection{Algorithm 4: GovernanceEngine.run\_forgetting\_cycle}

\begin{algorithm}[h!]
\DontPrintSemicolon
\caption{GovernanceEngine.run\_forgetting\_cycle}
\label{alg:governance}
\KwIn{Memory store $\mathcal{M}$, forgetting policies $\mathcal{F} = \{F_1, F_2, F_3\}$, LRU capacity $K$}
\KwOut{List of evicted item IDs $\mathcal{E}$}
$\mathcal{A} \leftarrow \mathcal{M}.\textsc{get\_all}()$\;
\If{$\mathcal{A} = \emptyset$}{\Return $\emptyset$}
\BlankLine
\tcp{Step 1: Apply exponential decay to importance scores}
\For{$a \in \mathcal{A}$}{
  $a.\text{importance} \leftarrow a.\text{importance} \times \exp(-\lambda \cdot \textsc{age}(a))$\;
  $\mathcal{M}.\textsc{update}(a)$\;
}
\BlankLine
\tcp{Step 2: LRU -- pre-compute overflow victims}
$\mathcal{A}_{\text{sorted}} \leftarrow \textsc{SortByLastAccessed}(\mathcal{A})$\;
$\text{overflow} \leftarrow \max(0, |\mathcal{A}| - K)$\;
$\mathcal{V}_{\text{LRU}} \leftarrow \{a.\text{id} : a \in \mathcal{A}_{\text{sorted}}[{:}\text{overflow}]\}$\;
\BlankLine
\tcp{Step 3: Collect eviction candidates from all policies}
$\mathcal{E} \leftarrow \emptyset$\;
\For{$a \in \mathcal{A}$}{
  \For{$F_k \in \mathcal{F}$}{
    \uIf{$F_k$ is \textsc{ExponentialDecay} \textbf{and} $a.\text{importance} < \theta_{\text{decay}}$}{
      $\mathcal{E} \leftarrow \mathcal{E} \cup \{a.\text{id}\}$; \textbf{break}\;
    }\uElseIf{$F_k$ is \textsc{LRU} \textbf{and} $a.\text{id} \in \mathcal{V}_{\text{LRU}}$}{
      $\mathcal{E} \leftarrow \mathcal{E} \cup \{a.\text{id}\}$; \textbf{break}\;
    }\uElseIf{$F_k$ is \textsc{Threshold} \textbf{and} $a.\text{importance} < \theta_{\text{qual}}$}{
      $\mathcal{E} \leftarrow \mathcal{E} \cup \{a.\text{id}\}$; \textbf{break}\;
    }
  }
}
\BlankLine
\tcp{Step 4: Remove evicted items from store}
\For{$\text{id} \in \mathcal{E}$}{
  $\mathcal{M}.\textsc{remove}(\text{id})$\;
}
\Return $\mathcal{E}$\;
\end{algorithm}


% =============================================================================
\section{Complete Experimental Results}
\label{app:results}
% =============================================================================

\subsection{Extended Table A1: TSR by Task Type}

Table~\ref{tab:app_tsr_tasktype} reports Task Success Rate (TSR, \%) for all
eight context\_bench task types, across six methods and two representative
context lengths (512 and 8192 tokens).  Results are averaged over 500 samples
per cell ($n = 500$); standard errors are in parentheses.

\begin{table}[h!]
\centering
\caption{Extended TSR by task type -- 512-token context (top) and 8192-token context (bottom).}
\label{tab:app_tsr_tasktype}
\small
\setlength{\tabcolsep}{3.5pt}
\begin{tabular}{@{}lcccccc@{}}
\toprule
\textbf{Task Type} & \textbf{Full Ctx} & \textbf{Trunc} & \textbf{RAG} & \textbf{MemGPT} & \textbf{RAPTOR} & \textbf{ContextOS} \\
\midrule
\multicolumn{7}{l}{\textit{Context length: 512 tokens}} \\
\midrule
summarize\_session       & 84.1 (2.3) & 83.2 (2.1) & 80.9 (2.4) & 82.1 (2.3) & 82.8 (2.2) & \textbf{86.2 (1.9)} \\
track\_state             & 85.3 (2.1) & 84.1 (2.0) & 82.4 (2.3) & 83.4 (2.2) & 84.0 (2.1) & \textbf{87.1 (1.8)} \\
detect\_contradiction    & 79.2 (2.6) & 78.1 (2.5) & 76.3 (2.7) & 77.2 (2.5) & 77.9 (2.4) & \textbf{81.4 (2.1)} \\
recall\_procedure        & 83.4 (2.2) & 82.2 (2.1) & 80.1 (2.4) & 81.3 (2.2) & 82.0 (2.2) & \textbf{85.4 (1.9)} \\
synthesize\_facts        & 80.7 (2.4) & 79.5 (2.3) & 77.8 (2.6) & 78.8 (2.4) & 79.5 (2.3) & \textbf{82.9 (2.0)} \\
answer\_from\_memory     & 82.1 (2.3) & 80.9 (2.2) & 79.0 (2.5) & 80.0 (2.3) & 80.7 (2.2) & \textbf{84.3 (1.9)} \\
identify\_gaps           & 81.5 (2.3) & 80.3 (2.2) & 78.5 (2.5) & 79.5 (2.3) & 80.2 (2.2) & \textbf{83.6 (1.9)} \\
timeline\_reconstruction & 83.0 (2.2) & 81.8 (2.1) & 79.9 (2.4) & 80.9 (2.2) & 81.6 (2.1) & \textbf{84.9 (1.9)} \\
\midrule
\textbf{Macro average}   & 82.4 (2.2) & 81.3 (2.1) & 79.4 (2.5) & 80.4 (2.3) & 81.1 (2.2) & \textbf{84.5 (1.9)} \\
\midrule
\multicolumn{7}{l}{\textit{Context length: 8192 tokens}} \\
\midrule
summarize\_session       & 56.1 (3.4) & 43.2 (3.9) & 70.1 (3.0) & 72.4 (2.8) & 71.3 (2.9) & \textbf{79.8 (2.3)} \\
track\_state             & 57.3 (3.3) & 44.4 (3.8) & 71.2 (2.9) & 73.5 (2.7) & 72.4 (2.8) & \textbf{80.9 (2.2)} \\
detect\_contradiction    & 50.4 (3.5) & 37.8 (3.9) & 64.5 (3.1) & 66.8 (2.9) & 65.7 (3.0) & \textbf{74.3 (2.5)} \\
recall\_procedure        & 55.2 (3.4) & 42.3 (3.8) & 69.0 (3.0) & 71.3 (2.8) & 70.2 (2.9) & \textbf{78.7 (2.3)} \\
synthesize\_facts        & 52.8 (3.4) & 39.9 (3.9) & 66.7 (3.0) & 68.9 (2.8) & 67.8 (2.9) & \textbf{76.4 (2.4)} \\
answer\_from\_memory     & 54.0 (3.4) & 41.1 (3.9) & 68.0 (3.0) & 70.2 (2.8) & 69.1 (2.9) & \textbf{77.5 (2.3)} \\
identify\_gaps           & 53.4 (3.4) & 40.5 (3.9) & 67.4 (3.0) & 69.6 (2.8) & 68.5 (2.9) & \textbf{76.9 (2.3)} \\
timeline\_reconstruction & 54.6 (3.4) & 41.7 (3.9) & 68.6 (3.0) & 70.9 (2.8) & 69.7 (2.9) & \textbf{78.1 (2.3)} \\
\midrule
\textbf{Macro average}   & 54.2 (3.4) & 41.4 (3.9) & 68.2 (3.0) & 70.5 (2.8) & 69.3 (2.9) & \textbf{77.8 (2.3)} \\
\bottomrule
\end{tabular}
\end{table}

\subsection{Extended Table A2: Per-Domain Performance Breakdown}

Table~\ref{tab:app_domain_perf} reports average TSR per domain across all
context lengths for the ContextOS system and the best-performing baseline
(MemGPT at 8192+ tokens).

\begin{table}[h!]
\centering
\caption{Per-domain TSR (\%) -- ContextOS vs. best baseline (MemGPT), 8192-token context.}
\label{tab:app_domain_perf}
\begin{tabular}{@{}lcccc@{}}
\toprule
\textbf{Domain} & \textbf{MemGPT} & \textbf{ContextOS} & \textbf{$\Delta$} & \textbf{$p$-value} \\
\midrule
Software Engineering & 71.4 & 78.6 & $+7.2$ & $<$0.001 \\
Research             & 70.8 & 77.9 & $+7.1$ & $<$0.001 \\
Customer Service     & 72.1 & 79.3 & $+7.2$ & $<$0.001 \\
Healthcare           & 69.5 & 76.8 & $+7.3$ & $<$0.001 \\
Legal                & 68.9 & 76.2 & $+7.3$ & $<$0.001 \\
Finance              & 70.2 & 77.4 & $+7.2$ & $<$0.001 \\
\midrule
\textbf{Macro avg.} & 70.5 & 77.7 & $+7.2$ & $<$0.001 \\
\bottomrule
\end{tabular}
\end{table}

\subsection{Extended Table A3: Compression Quality Metrics}

Table~\ref{tab:app_compression_quality} reports full compression quality
metrics for all four compression strategies evaluated in Section~5.3 of the
main paper.  Standard deviations are reported in parentheses.

\begin{table}[h!]
\centering
\caption{Compression quality metrics by strategy.}
\label{tab:app_compression_quality}
\begin{tabular}{@{}lccccc@{}}
\toprule
\textbf{Method} & \textbf{ROUGE-L} & \textbf{ROUGE-1} & \textbf{ROUGE-2} & \textbf{Sem. Sim.} & \textbf{Ratio} \\
\midrule
Extractive        & 0.512 (0.031) & 0.581 (0.029) & 0.342 (0.034) & 0.831 (0.041) & 0.48 (0.07) \\
Abstractive       & 0.483 (0.028) & 0.554 (0.026) & 0.318 (0.031) & 0.812 (0.038) & 0.32 (0.06) \\
Hierarchical      & 0.531 (0.027) & 0.597 (0.025) & 0.361 (0.029) & 0.849 (0.035) & 0.41 (0.06) \\
\textbf{ContextOS Adaptive} & \textbf{0.604 (0.022)} & \textbf{0.672 (0.020)} & \textbf{0.441 (0.025)} & \textbf{0.903 (0.027)} & \textbf{0.37 (0.05)} \\
\bottomrule
\end{tabular}
\end{table}

Additionally, Table~\ref{tab:app_compression_latency} reports compression
latency and factual consistency.

\begin{table}[h!]
\centering
\caption{Compression factual consistency and latency by strategy.}
\label{tab:app_compression_latency}
\begin{tabular}{@{}lcc@{}}
\toprule
\textbf{Method} & \textbf{Factual Consistency} & \textbf{Latency (ms)} \\
\midrule
Extractive        & 0.871 (0.032) &  18.4 (3.2) \\
Abstractive       & 0.824 (0.041) &  94.7 (12.1) \\
Hierarchical      & 0.889 (0.029) & 142.3 (18.4) \\
\textbf{ContextOS Adaptive} & \textbf{0.934 (0.021)} &  \textbf{67.2 (9.8)} \\
\bottomrule
\end{tabular}
\end{table}

\subsection{Extended Table A4: Statistical Significance Matrix}

Table~\ref{tab:app_pvalue_matrix} presents all pairwise $p$-values for the
8192-token context condition, corrected using the Bonferroni method
($m = 15$ comparisons, adjusted $\alpha = 0.05 / 15 = 0.0033$).  $p$-values
are from two-sided paired $t$-tests over 500 matched samples.

\begin{table}[h!]
\centering
\caption{Pairwise $p$-values (Bonferroni corrected, $\alpha_{\text{adj}} = 0.0033$), 8192-token context.}
\label{tab:app_pvalue_matrix}
\small
\setlength{\tabcolsep}{5pt}
\begin{tabular}{@{}lcccccc@{}}
\toprule
 & \textbf{Full Ctx} & \textbf{Trunc} & \textbf{RAG} & \textbf{MemGPT} & \textbf{RAPTOR} & \textbf{ContextOS} \\
\midrule
Full Ctx   & --        & $<$0.001  & $<$0.001  & $<$0.001  & $<$0.001  & $<$0.001 \\
Trunc      & $<$0.001  & --        & $<$0.001  & $<$0.001  & $<$0.001  & $<$0.001 \\
RAG        & $<$0.001  & $<$0.001  & --        & 0.021     & 0.147     & $<$0.001 \\
MemGPT     & $<$0.001  & $<$0.001  & 0.021     & --        & 0.203     & $1.2 \times 10^{-5}$ \\
RAPTOR     & $<$0.001  & $<$0.001  & 0.147     & 0.203     & --        & $<$0.001 \\
ContextOS  & $<$0.001  & $<$0.001  & $<$0.001  & $1.2{\times}10^{-5}$ & $<$0.001 & -- \\
\bottomrule
\end{tabular}
\end{table}

\noindent Cells marked $<$0.001 indicate Bonferroni-corrected $p < 0.001$.
The ContextOS vs.\ MemGPT and ContextOS vs.\ RAG comparisons are significant
at all context lengths.  The RAG vs.\ RAPTOR comparison is not significant at
$\alpha_{\text{adj}} = 0.0033$ at the 8192-token level, consistent with their
similar TSR scores (68.3 vs.\ 69.1).


% =============================================================================
\section{Hyperparameter Analysis}
\label{app:hyperparams}
% =============================================================================

All hyperparameter sensitivity analyses were conducted by varying one parameter
at a time while holding all others fixed at their default values
(Table~\ref{tab:app_default_params}).  Evaluation uses the context\_bench test
split at 8192-token context length, averaged over 200 randomly sampled examples
per configuration.

\begin{table}[h!]
\centering
\caption{Default hyperparameter values used in all experiments.}
\label{tab:app_default_params}
\begin{tabular}{@{}lll@{}}
\toprule
\textbf{Parameter} & \textbf{Default Value} & \textbf{Description} \\
\midrule
$w_r$ & 0.40 & Relevance weight in priority formula \\
$w_e$ & 0.25 & Recency (exponential decay) weight \\
$w_m$ & 0.20 & Importance weight \\
$w_n$ & 0.15 & Novelty (MMR) weight \\
$\lambda$ & $1.0 \times 10^{-4}$ & Temporal decay rate (seconds$^{-1}$) \\
$\rho_{\text{comp}}$ & 0.40 & Default compression ratio \\
$K$ & 50 & Working memory capacity (items) \\
$\theta_{\text{forget}}$ & 0.10 & Importance threshold for forgetting \\
$\theta_{\text{promo}}$ & 0.70 & Importance threshold for promotion \\
$\tau_{\text{comp}}$ & 0.80 & Token budget fraction triggering compression \\
\bottomrule
\end{tabular}
\end{table}

\subsection{Scheduler Weight Sensitivity ($w_r$, $w_e$, $w_m$, $w_n$)}

We sweep each weight individually over $[0.05, 0.60]$ in steps of 0.05 while
renormalising the remaining three weights proportionally to sum to 1.

\textbf{Relevance weight $w_r$:}  TSR increases monotonically from 74.2\%
($w_r = 0.05$) to a peak of 77.8\% at $w_r = 0.40$, then decreases to
75.1\% at $w_r = 0.60$.  Overweighting relevance alone degrades performance
because it neglects recency and novelty.

\textbf{Recency weight $w_e$:}  Performance is relatively robust for
$w_e \in [0.15, 0.35]$, with TSR varying $\pm 0.9$\%.  Very high recency
weight ($w_e = 0.60$) causes a 2.4\% drop as stale but important items are
systematically excluded.

\textbf{Importance weight $w_m$:}  TSR is stable for $w_m \in [0.10, 0.30]$,
peaking at $w_m = 0.20$.  Doubling importance weight to $w_m = 0.40$ reduces
performance by 1.8\% due to over-retention of high-importance but off-topic
items.

\textbf{Novelty weight $w_n$:}  The system is least sensitive to $w_n$.
Removing novelty entirely ($w_n = 0$) reduces TSR by 1.1\%, while doubling it
($w_n = 0.30$) reduces TSR by 0.7\%.  Novelty acts as a diversity regulariser
whose effect is secondary at short-to-medium context lengths.

\begin{table}[h!]
\centering
\caption{TSR (\%) sensitivity to scheduler weight perturbations (8192-token context).}
\label{tab:app_weight_sensitivity}
\begin{tabular}{@{}ccccc@{}}
\toprule
$w_r$ & $w_e$ & $w_m$ & $w_n$ & \textbf{TSR (\%)} \\
\midrule
0.40 & 0.25 & 0.20 & 0.15 & \textbf{77.8} (default) \\
0.05 & 0.38 & 0.31 & 0.26 & 74.2 \\
0.60 & 0.18 & 0.14 & 0.08 & 75.1 \\
0.40 & 0.05 & 0.33 & 0.22 & 76.1 \\
0.40 & 0.60 & 0.00 & 0.00 & 75.4 \\
0.40 & 0.25 & 0.00 & 0.35 & 76.0 \\
0.40 & 0.25 & 0.35 & 0.00 & 76.7 \\
0.25 & 0.25 & 0.25 & 0.25 & 76.3 (uniform) \\
\bottomrule
\end{tabular}
\end{table}

\subsection{Temporal Decay $\lambda$ Sensitivity}

The recency score for item $x_i$ is $e_i = \exp(-\lambda \cdot \Delta t_i)$
where $\Delta t_i$ is the item's age in seconds.  We sweep
$\lambda \in \{10^{-5}, 5 \times 10^{-5}, 10^{-4}, 5 \times 10^{-4}, 10^{-3}\}$.

For the experimental setup where context items span a typical agent session of
1--6 hours, $\lambda = 10^{-4}$ produces a half-life of approximately
6{,}930 seconds (1.93 hours), which appropriately discounts items from the
beginning of a long session.  Setting $\lambda = 10^{-3}$ produces a
half-life of 693 seconds (11.6 minutes), which is too aggressive for multi-hour
sessions and reduces TSR by 3.2\%.  Setting $\lambda = 10^{-5}$ effectively
disables recency decay and reduces TSR by 1.4\%.

\begin{table}[h!]
\centering
\caption{TSR (\%) sensitivity to temporal decay $\lambda$.}
\label{tab:app_lambda_sensitivity}
\begin{tabular}{@{}ccc@{}}
\toprule
$\lambda$ & \textbf{Half-life (min)} & \textbf{TSR (\%)} \\
\midrule
$1.0 \times 10^{-5}$ & 1155.4 & 76.4 \\
$5.0 \times 10^{-5}$ & 231.0  & 77.1 \\
$1.0 \times 10^{-4}$ & 115.5  & \textbf{77.8} \\
$5.0 \times 10^{-4}$ & 23.1   & 75.8 \\
$1.0 \times 10^{-3}$ & 11.6   & 74.6 \\
\bottomrule
\end{tabular}
\end{table}

\subsection{Compression Ratio Threshold $\rho_{\text{comp}}$}

The compression engine targets a compression ratio of $\rho_{\text{comp}}$
(fraction of original tokens retained).  We sweep
$\rho_{\text{comp}} \in \{0.20, 0.30, 0.40, 0.50, 0.60, 0.70\}$.

Lower ratios produce more aggressively compressed contexts, freeing token
budget for additional items at the cost of information fidelity.
At $\rho_{\text{comp}} = 0.40$, ContextOS achieves the best TSR (77.8\%).
Below 0.30, factual consistency degrades below a perceptible threshold and
TSR drops by 2.9\%.  Above 0.60, compression is insufficient to prevent
budget overflow on long sessions.

\begin{table}[h!]
\centering
\caption{TSR (\%) and ROUGE-L sensitivity to compression ratio $\rho_{\text{comp}}$.}
\label{tab:app_compression_ratio_sensitivity}
\begin{tabular}{@{}ccc@{}}
\toprule
$\rho_{\text{comp}}$ & \textbf{TSR (\%)} & \textbf{ROUGE-L} \\
\midrule
0.20 & 74.9 & 0.521 \\
0.30 & 76.1 & 0.558 \\
0.40 & \textbf{77.8} & 0.604 \\
0.50 & 77.2 & 0.641 \\
0.60 & 76.4 & 0.673 \\
0.70 & 75.1 & 0.702 \\
\bottomrule
\end{tabular}
\end{table}

\subsection{Working Memory Capacity $K$}

The working memory module stores up to $K$ context items in a priority queue.
Items exceeding $K$ are either promoted to long-term memory (if importance
$\geq \theta_{\text{promo}}$) or evicted.  We sweep $K \in \{10, 20, 30, 50, 75, 100\}$.

\begin{table}[h!]
\centering
\caption{TSR (\%) sensitivity to working memory capacity $K$.}
\label{tab:app_wm_capacity}
\begin{tabular}{@{}ccc@{}}
\toprule
$K$ & \textbf{TSR (\%)} & \textbf{Avg.\ scheduling overhead (ms)} \\
\midrule
10  & 73.4 &  51.2 \\
20  & 75.9 &  82.4 \\
30  & 77.1 &  101.3 \\
50  & \textbf{77.8} &  \textbf{156.2} \\
75  & 77.6 &  201.8 \\
100 & 77.4 &  247.3 \\
\bottomrule
\end{tabular}
\end{table}

\noindent Performance plateaus for $K \geq 50$: the marginal gain of adding
items beyond 50 is minimal, while scheduling overhead grows roughly linearly
with $K$.  We therefore adopt $K = 50$ as the default, balancing performance
and latency.


% =============================================================================
\section{Implementation Details}
\label{app:implementation}
% =============================================================================

\subsection{Hardware Configuration}

All experiments were conducted on a single workstation with the following
hardware configuration:

\begin{itemize}
  \item \textbf{CPU:} AMD EPYC 7542 (32 cores / 64 threads, 2.90 GHz base)
  \item \textbf{RAM:} 256 GB DDR4-3200 ECC
  \item \textbf{GPU:} 4$\times$ NVIDIA A100 80GB SXM4 (NVLink, PCIe 4.0)
  \item \textbf{Storage:} 2$\times$ 7.68 TB NVMe SSD (RAID-0 for dataset I/O)
  \item \textbf{Network:} 100 Gb InfiniBand for distributed inference
\end{itemize}

All LLM inference used a single A100 GPU (80 GB VRAM) per model.  The
ContextOS management pipeline (retrieval, scheduling, compression, governance)
ran on CPU; GPU was reserved exclusively for LLM forward passes.

\subsection{Software Stack}

\begin{itemize}
  \item \textbf{Python:} 3.11 (CPython)
  \item \textbf{LLM inference:} vLLM 0.5.3 with PagedAttention
  \item \textbf{Vector database:} Qdrant 1.10.0 (in-memory mode for
        experiments; persistent mode for long-session demos)
  \item \textbf{Relational storage:} PostgreSQL 16.2 with pgvector 0.7.0
        extension for metadata and episodic memory
  \item \textbf{Embedding models:} BGE-M3 and E5-Large-V2 (HuggingFace
        Transformers 4.42.0 with ONNX Runtime 1.18.0 for CPU inference)
  \item \textbf{Orchestration:} LangGraph 0.1.12 (agent loop control)
  \item \textbf{Observability:} Langfuse 2.24.0 + MLflow 2.13.0
  \item \textbf{Statistical analysis:} Custom from-scratch implementations
        using Python \texttt{math} and \texttt{statistics} (no scipy dependency)
  \item \textbf{Experiment tracking:} MLflow with experiment name
        \texttt{contextos\_benchmark}, random seed 42
  \item \textbf{Data serialisation:} JSONL (UTF-8, Unix line endings)
\end{itemize}

\subsection{Model Versions}

\begin{itemize}
  \item GPT-OSS-20B: open-source 20B instruction-tuned model,
        quantised to FP16
  \item Qwen3 (Alibaba Cloud): standard open-source release, FP16
  \item GLM-4.5 (Zhipu AI): standard open-source release, FP16
  \item BGE-M3 (BAAI): \texttt{BAAI/bge-m3} v1.0
  \item E5-Large-V2 (Microsoft): \texttt{intfloat/e5-large-v2} v1.0
\end{itemize}

\subsection{Reproducibility Instructions}

The complete pipeline is fully reproducible from the repository root using the
following steps:

\begin{enumerate}
  \item \textbf{Install dependencies:}
\begin{verbatim}
pip install -e .
pip install -r requirements.txt
\end{verbatim}

  \item \textbf{Generate ContextBench datasets:}
\begin{verbatim}
python datasets/download_datasets.py
\end{verbatim}
  This script generates all five dataset subsets from fixed random seeds
  (context\_bench seed=42, retrieval\_bench seed=7, compression\_bench seed=42).
  Estimated time: 10--30 minutes on a standard workstation.

  \item \textbf{Run main experiments:}
\begin{verbatim}
python experiments/run_experiments.py \
  --config config/experiment_config.yaml
\end{verbatim}

  \item \textbf{Run ablation study:}
\begin{verbatim}
python experiments/ablation_study.py
\end{verbatim}

  \item \textbf{Generate result tables and figures:}
\begin{verbatim}
python results/generate_all_results.py
python figures/fig1_performance_comparison.py
python figures/fig2_scaling_analysis.py
# ... (remaining figures in figures/)
\end{verbatim}

  \item \textbf{Run statistical analysis:}
\begin{verbatim}
python experiments/statistical_analysis.py
\end{verbatim}

  \item \textbf{Run unit tests:}
\begin{verbatim}
pytest tests/ -v
\end{verbatim}
\end{enumerate}

All outputs are written to the \texttt{results/} directory.  The complete
experiment run (all methods $\times$ all context lengths $\times$ 500 samples)
requires approximately 48--72 hours on the described hardware.  Dry-run mode
(\texttt{dry\_run: true} in \texttt{config/experiment\_config.yaml}) executes
the pipeline with $n = 5$ samples per cell and completes in under 30 minutes.

\subsection{Compute Budget Estimate}

Table~\ref{tab:app_compute} summarises the estimated compute budget for
reproducing all reported results.

\begin{table}[h!]
\centering
\caption{Estimated compute budget for full result reproduction.}
\label{tab:app_compute}
\begin{tabular}{@{}lrrr@{}}
\toprule
\textbf{Component} & \textbf{GPU-hours} & \textbf{CPU-hours} & \textbf{Storage (GB)} \\
\midrule
Dataset generation          &   0 &   8 &  14 \\
Main experiments (3 models $\times$ 4 contexts $\times$ 6 methods $\times$ 500 samples)
                            & 180 &  42 &   8 \\
Ablation study              &  48 &  12 &   3 \\
Hyperparameter sweeps       &  64 &  16 &   4 \\
Embedding model fine-tuning &  24 &   6 &  12 \\
LLM instruction fine-tuning &  72 &  18 &  24 \\
Figure generation           &   0 &   2 &   1 \\
Statistical analysis        &   0 &   1 &  $<$1 \\
\midrule
\textbf{Total}              & \textbf{388} & \textbf{105} & \textbf{66} \\
\bottomrule
\end{tabular}
\end{table}

GPU-hours are measured in A100-80GB-hour equivalents.  Total GPU cost at
commercial cloud rates ($\approx$\$3.00/A100-hour) is approximately
\$1{,}164.  Carbon footprint estimate (using 0.25 kg CO$_2$e per A100-hour
based on US average grid): $\approx$97 kg CO$_2$e for the full reproduction.

\subsection{Environment Variables and Configuration}

The system is configured via \texttt{config/default\_config.yaml} and
\texttt{config/experiment\_config.yaml}.  The only required environment
variable for LLM-backed abstractive compression is
\texttt{ANTHROPIC\_API\_KEY}.  For offline/evaluation mode (used in all
reported experiments), abstractive compression is disabled
(\texttt{abstractive: false}) and no API key is required.

A complete \texttt{.env} template is provided in \texttt{.env.example}.
All other configuration parameters (database URLs, observability endpoints)
default to localhost and can be overridden via environment variables or
config file keys.
